\documentclass[12pt,a4paper]{article}

\usepackage[T2A]{fontenc}
\usepackage[utf8]{inputenc}
\usepackage[russian]{babel}
\usepackage{amsmath,amssymb,amsthm,mathtools}
\usepackage{geometry}
\usepackage{microtype}
\usepackage{enumitem}
\usepackage{xcolor}
\usepackage{tikz}
\usepackage{hyperref}

\geometry{left=25mm,right=20mm,top=22mm,bottom=24mm}
\hypersetup{
  colorlinks=true,
  linkcolor=blue!45!black,
  urlcolor=blue!55!black
}
\setlist{nosep}
\setlength{\parindent}{1.25em}
\setlength{\parskip}{0.15em}

\newtheorem{theorem}{Теорема}
\newtheorem{proposition}{Утверждение}
\newtheorem{lemma}{Лемма}
\theoremstyle{definition}
\newtheorem{definition}{Определение}
\theoremstyle{remark}
\newtheorem*{remark}{Замечание}

\DeclareMathOperator{\Sp}{Sp}
\DeclareMathOperator{\diag}{diag}
\DeclareMathOperator{\cond}{cond}
\newcommand{\CC}{\mathbb{C}}
\newcommand{\RR}{\mathbb{R}}
\newcommand{\ZZ}{\mathbb{Z}}
\newcommand{\eps}{\varepsilon}
\newcommand{\norm}[1]{\left\lVert #1\right\rVert}

\DeclareMathOperator{\ind}{ind}
\newcommand{\Conv}{\mathcal C}

\title{Индекс и факторизация в алгебре Винера\\Ряд Неймана в банаховой алгебре}
\author{Конспект лекции}
\date{17 октября 2026 года}

\begin{document}
\maketitle

\section{Индекс и приращение логарифма}

Пусть $\Gamma=\{t\in\CC:|t|=1\}$ и $a\colon\Gamma\to\CC$ ---
непрерывная функция, не обращающаяся в нуль. При положительном обходе
окружности положим
\[
  g(\varphi)=a(e^{i\varphi}),\qquad 0\leqslant\varphi\leqslant2\pi.
\]
На предыдущей лекции индекс был определён через непрерывный аргумент
$\theta(\varphi)$ функции $g$:
\[
  g(\varphi)=|g(\varphi)|e^{i\theta(\varphi)},\qquad
  \ind a=\frac{\theta(2\pi)-\theta(0)}{2\pi}.
\]
Индекс --- целое число оборотов образа вокруг нуля с учётом направления.

\subsection{Непрерывный логарифм вдоль пути}

Для ненулевого комплексного числа логарифм имеет значения
\[
  \log z=\ln|z|+i(\arg z+2\pi k),\qquad k\in\ZZ.
\]
Например, все значения логарифма числа $-1$ равны
\[
  i(\pi+2\pi k),\qquad
  e^{i(\pi+2\pi k)}=-1.
\]
Вдоль отрезка параметра выберем непрерывный логарифм
\[
  L(\varphi)=\ln|g(\varphi)|+i\theta(\varphi).
\]
Он удовлетворяет $e^{L(\varphi)}=g(\varphi)$, но не обязан быть
периодическим: при возвращении $g$ в начальную точку значение $L$
может измениться на целое кратное $2\pi i$.

\begin{theorem}
Индекс выражается через приращение непрерывного логарифма вдоль пути:
\[
  \boxed{\ind a=\frac{L(2\pi)-L(0)}{2\pi i}}.
\]
\end{theorem}

\begin{proof}
Так как $g(2\pi)=g(0)$, модули в концах совпадают. Поэтому
\begin{align*}
  L(2\pi)-L(0)
  &=\ln|g(2\pi)|-\ln|g(0)|
    +i\bigl(\theta(2\pi)-\theta(0)\bigr)\\
  &=2\pi i\,\ind a.
\end{align*}
\end{proof}

\begin{remark}
Запись из рукописи $[\ln a(e^{i\varphi})]_0^{2\pi}$ понимается
именно как приращение непрерывно выбранного логарифма вдоль пути.
Подстановка двух главных значений логарифма в совпадающих концах
всегда дала бы нуль. Для этой формулы достаточно непрерывности;
дифференцируемость не требуется.
\end{remark}

\subsection{Индекс произведения}

\begin{theorem}
Для непрерывных ненулевых функций $a,b$ на $\Gamma$ справедливо
\[
  \boxed{\ind(ab)=\ind a+\ind b}.
\]
\end{theorem}

\begin{proof}
Если $\theta_a,\theta_b$ --- непрерывные аргументы соответствующих
путей, то $\theta_a+\theta_b$ --- непрерывный аргумент произведения.
Его полное приращение равно сумме приращений.
\end{proof}

Отсюда получаем
\[
  \ind(1/a)=-\ind a,\qquad
  \ind(a/b)=\ind a-\ind b.
\]
У ненулевой постоянной индекс равен нулю, а $\ind(t^m)=m$ для
любого $m\in\ZZ$.

\section{Индекс рациональной функции}

Пусть
\[
  a(t)=\frac{P(t)}{Q(t)},\qquad P(t)\ne0,\ Q(t)\ne0
  \quad(t\in\Gamma),
\]
где $P,Q$ --- ненулевые многочлены. Корни считаем с кратностями.

\begin{theorem}
Если $N$ --- число корней $P$ в $|z|<1$, а $M$ --- число корней
$Q$ в этом круге, то
\[
  \boxed{\ind a=N-M}.
\]
Для несократимой дроби это число нулей минус число полюсов внутри круга.
\end{theorem}

\begin{proof}
Разложим многочлены на линейные множители и применим аддитивность
индекса. Достаточно выяснить индекс $t-\alpha$.

Если $|\alpha|<1$, то
\[
  t-\alpha=t\left(1-\frac{\alpha}{t}\right).
\]
При $|t|=1$ второй множитель лежит в круге с центром $1$ и радиусом
$|\alpha|<1$, то есть в правой полуплоскости. Там аргумент можно
непрерывно выбрать в $(-\pi/2,\pi/2)$, и после полного обхода он
возвращается к начальному значению. Следовательно,
$\ind(1-\alpha/t)=0$ и $\ind(t-\alpha)=1$.

Если $|\alpha|>1$, то
\[
  t-\alpha=-\alpha\left(1-\frac{t}{\alpha}\right).
\]
По той же причине второй множитель имеет индекс нуль. Постоянный
множитель также имеет индекс нуль, поэтому $\ind(t-\alpha)=0$.
Каждый внутренний корень числителя даёт $+1$, а каждый внутренний
корень знаменателя --- $-1$.
\end{proof}

\begin{remark}
Общие корни $P,Q$ можно сократить: их вклады в разность $N-M$
взаимно уничтожаются. Корни на $\Gamma$ исключены условиями теоремы.
Доказательство через линейные множители разворачивает утверждение,
которое в рукописи дано без доказательства.
\end{remark}

\section{Алгебра Винера и её подалгебры}

Напомним определение:
\[
  W=\left\{a(t)=\sum_{k\in\ZZ}a_kt^k:
       \sum_{k\in\ZZ}|a_k|<\infty\right\},\qquad
  \norm{a}_W=\sum_{k\in\ZZ}|a_k|.
\]
Это коммутативная банахова алгебра относительно поточечного умножения.
Её единица --- постоянная функция $1$, причём $\norm{1}_W=1$.

\begin{theorem}[Винера]
Если $a\in W$ и $a(t)\ne0$ для всех $t\in\Gamma$, то $1/a\in W$.
Иначе говоря, обратная функция также имеет абсолютно сходящийся
ряд Фурье.
\end{theorem}

В рукописи эта теорема сформулирована без доказательства; здесь она
также используется как известный результат. Обратное утверждение
непосредственно: если $a$ обратима в алгебре функций, она не может
обращаться в нуль на окружности.

\subsection{Неотрицательные и неположительные степени}

\begin{definition}
Обозначим
\begin{align*}
  W_+&=\left\{a\in W:a_k=0\text{ при }k<0\right\},\\
  W_-&=\left\{a\in W:a_k=0\text{ при }k>0\right\}.
\end{align*}
Таким образом, в $W_+$ разрешены только неотрицательные степени $t$,
а в $W_-$ --- только неположительные.
\end{definition}

Оба множества --- замкнутые подалгебры $W$ с единицей. Действительно,
сложение и умножение сохраняют соответствующие знаки степеней, а
сходимость в норме $W$ сохраняет нулевые коэффициенты. Кроме того,
\[
  a(t)\in W_+\quad\Longleftrightarrow\quad a(1/t)\in W_-,
  \qquad W_+\cap W_-=\CC\cdot1.
\]
Последнее равенство следует из однозначности коэффициентов Фурье:
в пересечении могут остаться только постоянные функции.

Ряд элемента $W_+$ задаёт аналитическую функцию внутри единичного
круга и непрерывную на его замыкании. Ряд элемента $W_-$ задаёт
аналитическую функцию при $|z|>1$, непрерывную вплоть до окружности,
с конечным значением в бесконечности, равным $a_0$.

\subsection{Две леммы о линейных множителях}

\begin{lemma}
Если $|\alpha|>1$, то $\alpha-t$ и $(\alpha-t)^{-1}$ принадлежат $W_+$.
\end{lemma}

\begin{proof}
Первое утверждение очевидно. Для второго используем геометрический ряд:
\[
  \frac1{\alpha-t}=\frac1\alpha\frac1{1-t/\alpha}
  =\sum_{k=0}^{\infty}\frac{t^k}{\alpha^{k+1}}.
\]
Его коэффициенты абсолютно суммируемы:
\[
  \sum_{k=0}^{\infty}\frac1{|\alpha|^{k+1}}
  =\frac1{|\alpha|-1}<\infty.
\]
\end{proof}

\begin{lemma}
Если $|\beta|<1$, то $1-\beta/t$ и $(1-\beta/t)^{-1}$ принадлежат $W_-$.
\end{lemma}

\begin{proof}
На окружности $|\beta/t|=|\beta|<1$, поэтому
\[
  \frac1{1-\beta/t}=\sum_{k=0}^{\infty}\beta^kt^{-k},\qquad
  \sum_{k=0}^{\infty}|\beta|^k=\frac1{1-|\beta|}<\infty.
\]
Все степени неположительны.
\end{proof}

\begin{remark}
Принадлежность $a$ подалгебре $W_+$ сама по себе не означает
$1/a\in W_+$. Например, $t\in W_+$, но $1/t\in W_-$ и не принадлежит
$W_+$. Поэтому в факторизации отдельно требуется принадлежность
обратных множителей тем же подалгебрам.
\end{remark}

\section{Факторизация рациональной функции}

\begin{theorem}
Пусть $a=P/Q$ и оба многочлена не имеют корней на $\Gamma$.
Тогда существуют функции $a_-,a_+$ такие, что
\[
  \boxed{a(t)=a_-(t)t^{\kappa}a_+(t)},\qquad
  a_\pm,a_\pm^{-1}\in W_\pm,\qquad \kappa=\ind a.
\]
\end{theorem}

\begin{proof}
Разделим корни каждого многочлена на внутренние и внешние:
\begin{align*}
  P(t)&=p\prod_{j=1}^{N}(t-\alpha_j)
          \prod_{j=1}^{r}(t-\gamma_j),\\
  Q(t)&=q\prod_{j=1}^{M}(t-\beta_j)
          \prod_{j=1}^{s}(t-\delta_j),
\end{align*}
где $p,q\ne0$,
\[
  |\alpha_j|,|\beta_j|<1,\qquad |\gamma_j|,|\delta_j|>1.
\]
Все корни перечислены с кратностями; пустое произведение равно $1$.
Для внутренних корней вынесем $t$ из каждого линейного множителя.
Получим
\[
  a(t)=
  \underbrace{\frac{\prod_{j=1}^{N}(1-\alpha_j/t)}
                    {\prod_{j=1}^{M}(1-\beta_j/t)}}_{a_-(t)}
  t^{N-M}
  \underbrace{\frac pq
       \frac{\prod_{j=1}^{r}(t-\gamma_j)}
            {\prod_{j=1}^{s}(t-\delta_j)}}_{a_+(t)}.
\]
По двум леммам и замкнутости подалгебр относительно умножения
$a_-,a_-^{-1}\in W_-$ и $a_+,a_+^{-1}\in W_+$.
Из формулы индекса рациональной функции $N-M=\ind a$.
\end{proof}

\subsection{Единственность с точностью до постоянного множителя}

При фиксированной степени $\kappa$ пусть имеются две такие факторизации:
\[
  a_-t^\kappa a_+=\widetilde a_-t^\kappa\widetilde a_+.
\]
Тогда
\[
  \widetilde a_-^{-1}a_-
  =\widetilde a_+a_+^{-1}\in W_-\cap W_+.
\]
Обе части равны ненулевой постоянной. Следовательно, множители
определены с точностью до переноса постоянного коэффициента из одного
в другой. Условие $a_-(\infty)=1$ устраняет эту свободу.
В построении выше эта нормировка уже выполнена.

\begin{remark}
В рукописи единственность отмечена без доказательства. Здесь оно
добавлено через пересечение подалгебр. В этом разделе доказана
факторизация именно рациональной функции; общий случай произвольного
ненулевого элемента $W$ этим доказательством не охватывается.
\end{remark}

\subsection{Пример построения множителей}

В качестве дополнительного примера возьмём
\[
  a(t)=\frac{(t-\tfrac12)(t-2)}{t-\tfrac13}.
\]
Внутри круга по одному корню числителя и знаменателя, поэтому
$\kappa=1-1=0$. Нормированная факторизация имеет вид
\[
  a_-(t)=\frac{1-\tfrac1{2t}}{1-\tfrac1{3t}},\qquad
  a_+(t)=t-2,\qquad a(t)=a_-(t)a_+(t).
\]
Обратные множители принадлежат соответствующим подалгебрам по леммам.

\section{Ряд Неймана в банаховой алгебре}

Далее $\mathcal A$ --- банахова алгебра с единицей $e$, в которой
\[
  \norm{uv}\leqslant\norm{u}\norm{v},\qquad \norm{e}=1.
\]
Коммутативность в этом разделе не требуется. Нормировка единицы
явно добавлена: именно при ней оценки из рукописи имеют указанные
числовые константы.

\begin{theorem}
Если $u\in\mathcal A$ и $\norm{u}\leqslant q<1$, то $e-u$ обратим,
причём
\[
  \boxed{(e-u)^{-1}=\sum_{n=0}^{\infty}u^n},\qquad
  \boxed{\norm{(e-u)^{-1}}\leqslant\frac1{1-q}},
\]
где $u^0=e$.
\end{theorem}

\begin{proof}
Из субмультипликативности нормы следует $\norm{u^n}\leqslant q^n$.
Поэтому ряд абсолютно сходится в банаховой алгебре:
\[
  \sum_{n=0}^{\infty}\norm{u^n}\leqslant\sum_{n=0}^{\infty}q^n
  =\frac1{1-q}.
\]
Обозначим его сумму через $v$, а частичную сумму --- через
$S_N=e+u+\dots+u^N$. Конечное телескопическое тождество даёт
\[
  (e-u)S_N=S_N(e-u)=e-u^{N+1}.
\]
Так как $\norm{u^{N+1}}\leqslant q^{N+1}\to0$, переход к пределу
даёт $(e-u)v=v(e-u)=e$. Значит, $v=(e-u)^{-1}$, а оценка нормы
следует из оценки суммы ряда.
\end{proof}

Это аналог формулы геометрической прогрессии
$(1-z)^{-1}=\sum_{n\geqslant0}z^n$ при $|z|<1$.
В рукописи сокращение бесконечных рядов записано формально;
переход через $S_N$ обосновывает его.

Для приближённого вычисления полезна оценка остатка:
\[
  \left\lVert(e-u)^{-1}-\sum_{n=0}^{N}u^n\right\rVert
  \leqslant\frac{q^{N+1}}{1-q}.
\]
Это дополнительное следствие той же оценки геометрического ряда.

\section{Устойчивость обратимости при возмущении}

\begin{theorem}
Пусть $a\in\mathcal A$ обратим и
\[
  \norm{b-a}\leqslant\frac{q}{\norm{a^{-1}}},\qquad 0<q<1.
\]
Тогда $b$ также обратим и
\[
  \boxed{\norm{b^{-1}}\leqslant\frac{\norm{a^{-1}}}{1-q}}.
\]
\end{theorem}

\begin{proof}
Положим $d=a^{-1}(a-b)$. Тогда
\[
  b=a\bigl(e-a^{-1}(a-b)\bigr)=a(e-d),\qquad
  \norm{d}\leqslant\norm{a^{-1}}\norm{a-b}\leqslant q<1.
\]
По теореме о ряде Неймана $e-d$ обратим. Поэтому
\[
  b^{-1}=(e-d)^{-1}a^{-1},\qquad
  \norm{b^{-1}}\leqslant\norm{(e-d)^{-1}}\norm{a^{-1}}
  \leqslant\frac{\norm{a^{-1}}}{1-q}.
\]
Порядок множителей сохранён: алгебра может быть некоммутативной.
\end{proof}

В частности, множество обратимых элементов открыто: любое возмущение
с $\norm{b-a}<1/\norm{a^{-1}}$ сохраняет обратимость.
Дополнительно из тождества
\[
  b^{-1}-a^{-1}=b^{-1}(a-b)a^{-1}
\]
получаем оценку изменения обратного элемента
\[
  \norm{b^{-1}-a^{-1}}
  \leqslant\frac{\norm{a^{-1}}^2}{1-q}\norm{b-a}.
\]
Она показывает непрерывность операции обращения в каждой обратимой точке.

\section{Подстановка элемента алгебры в степенной ряд}

Пусть скалярный степенной ряд
\[
  f(z)=\sum_{n=0}^{\infty}c_nz^n
\]
имеет радиус сходимости $R>0$. Если $u\in\mathcal A$ и $\norm{u}<R$,
можно определить
\[
  \boxed{f(u)=\sum_{n=0}^{\infty}c_nu^n\in\mathcal A}.
\]
Постоянный член здесь означает $c_0e$. Корректность определения
следует из абсолютной сходимости:
\[
  \sum_{n=0}^{\infty}\norm{c_nu^n}
  \leqslant\sum_{n=0}^{\infty}|c_n|\norm{u}^n<\infty.
\]
Полнота $\mathcal A$ гарантирует существование суммы.

Для дальнейшего дополнения понадобятся два частных случая:
\begin{align*}
  \exp u&=\sum_{n=0}^{\infty}\frac{u^n}{n!}
       &&\text{для любого }u\in\mathcal A,\\
  \log(e-u)&=-\sum_{n=1}^{\infty}\frac{u^n}{n}
       &&\text{при }\norm{u}<1.
\end{align*}
Экспоненциальный ряд имеет бесконечный радиус сходимости, а
логарифмический --- радиус $1$. Эти примеры разворачивают определение
из рукописи и будут использованы для завершения последней теоремы.

\section{Факторизация функции, близкой к единице}

Вернёмся к алгебре Винера. Пусть $d\in W$ и
\[
  \norm{d}_W<1.
\]
Буква $d$ обозначает малое возмущение; она отделена от обозначений
корней многочленов предыдущих разделов.

\begin{theorem}
Существуют функции $a_-,a_+$ такие, что
\[
  \boxed{1-d(t)=a_-(t)a_+(t)},\qquad
  a_\pm,a_\pm^{-1}\in W_\pm.
\]
Степенной множитель здесь отсутствует: $\ind(1-d)=0$.
\end{theorem}

\subsection{Геометрическое пояснение из рукописи}

Для каждого $t\in\Gamma$ имеем
\[
  |d(t)|=\left|\sum_{k\in\ZZ}d_kt^k\right|
  \leqslant\sum_{k\in\ZZ}|d_k|=\norm{d}_W<1.
\]
Следовательно, образ функции $1-d$ лежит в замкнутом круге
с центром $1$ и радиусом $r=\norm{d}_W<1$. Этот круг целиком
расположен в правой полуплоскости и не содержит нуля. Непрерывный
аргумент возвращается к начальному значению, поэтому
\[
  \ind(1-d)=0.
\]

\begin{figure}[ht]
  \centering
  \begin{tikzpicture}[>=stealth,scale=2]
    \fill[blue!6] (1,0) circle (0.65);
    \draw[blue!60!black,thick] (1,0) circle (0.65);
    \draw[->,gray] (-0.35,0) -- (2,0) node[right] {$\operatorname{Re}z$};
    \draw[->,gray] (0,-0.9) -- (0,1) node[above] {$\operatorname{Im}z$};
    \fill (1,0) circle (0.025);
    \node[below left] at (0,0) {$0$};
    \node[below] at (1,0) {$1$};
    \draw[dashed] (1,0) -- node[above] {$r<1$} (1.46,0.46);
    \node[align=center] at (1,-1.05) {$(1-d)(\Gamma)\subset\{z:|z-1|\leqslant r\}$};
  \end{tikzpicture}
  \caption{Область, содержащая образ функции $1-d$}
\end{figure}

\subsection{Завершение доказательства: логарифм и разделение степеней}

\emph{Дополнение к рукописи.} На последней странице показана
геометрия образа и указан нулевой индекс; построение множителей
не записано. Завершим доказательство средствами степенных рядов.

Так как $\norm{d}_W<1$, в банаховой алгебре $W$ сходится ряд
\[
  h=\log(1-d)=-\sum_{n=1}^{\infty}\frac{d^n}{n}\in W,
  \qquad
  \norm{h}_W\leqslant\sum_{n=1}^{\infty}\frac{\norm{d}_W^n}{n}
  =-\ln(1-\norm{d}_W).
\]
В каждой точке $t\in\Gamma$ это обычный логарифмический ряд при
$|d(t)|<1$, поэтому $\exp h(t)=1-d(t)$.

Разделим ряд Фурье функции $h$ на отрицательные и неотрицательные степени:
\[
  h(t)=\sum_{k\in\ZZ}h_kt^k=h_-(t)+h_+(t),\qquad
  h_-\!:=\sum_{k<0}h_kt^k,\quad h_+\!:=\sum_{k\geqslant0}h_kt^k.
\]
Оба ряда абсолютно сходятся, причём $h_-\in W_-$ и $h_+\in W_+$.
Положим
\[
  a_-=\exp h_-,\qquad a_+=\exp h_+.
\]
Замкнутые подалгебры $W_\pm$ содержат единицу, поэтому
экспоненциальные ряды сходятся в соответствующих подалгебрах.
Точно так же
\[
  a_-^{-1}=\exp(-h_-)\in W_-,\qquad
  a_+^{-1}=\exp(-h_+)\in W_+.
\]
Поточечное умножение коммутативно; значит,
\[
  a_-a_+=\exp(h_-+h_+)=\exp h=1-d.
\]
Факторизация построена. Поскольку $h_-(\infty)=0$, выполнена
нормировка $a_-(\infty)=1$.

\begin{remark}
Условие $\norm{d}_W<1$ сильнее, чем $\max_{t\in\Gamma}|d(t)|<1$.
Именно оценка в норме Винера обеспечивает сходимость логарифмического
ряда в $W$. Геометрическое рассуждение само по себе устанавливает
ненулевость и индекс, но ещё не строит множители в нужных подалгебрах.
\end{remark}

\end{document}
